\section{机器学习与表示学习}
\subsection{前馈神经网络与课程事项}

前馈神经网络（feedforward neural network）通过一系列有方向的层，将输入变换为输出。神经元模型与激活函数决定各层的计算方式，梯度下降则通过调整参数完成训练。与并行程序设计的联系在于：网络的前向计算需要大量点积与矩阵运算，训练还需要计算、传播并汇总这些运算对应的梯度。

\begin{overviewbox}[title={课程事项：2026年9月24日}]
Lab 3 于当周周五截止；项目里程碑 1（Project milestone 1）即将发布。期中考试 1 的时间为 10 月 6 日 19:00--22:00。课后事项为提交 Lab 3、选择项目；阅读范围为教材第 4 版第 16 章，或第 5 版第 19 章。
\end{overviewbox}

\subsection{从输入输出规则到从样本中学习}

机器学习（machine learning）从数据中提取模式，用这些模式建立输入与输出之间的关系。对于难以完整写出程序规则的任务，可以提供带标签的数据，即已经配对的输入和期望输出，再迭代调整模型参数，使模型的答案接近期望答案。这个调整过程称为训练（training）；参数确定以后，由新输入计算输出的过程称为推断（inference）。分类、回归、语音转写和翻译都是机器学习任务。

数据驱动的方法依赖三个条件：有足够的计算能力反复训练模型，有能够表达目标关系的数据，以及存在难以手工完整描述规则的应用需求。GPU 与 CUDA 缩短了深度神经网络的训练与实验周期；传感器、网络服务和存储设施增加了可获得的数据；自动驾驶、智能设备、安全和医疗等需求则推动了相关应用。

计算速度提高以后，网络结构、特征选择、输入输出形式和高质量标注仍然需要设计。训练数据的作用是提供可学习的关系；数据是否覆盖了任务中的关键变化，直接影响模型有机会学到什么。

\subsection{穷举规则的规模与特征的作用}

\subsubsection{小规模布尔函数可以由真值表完整指定}
对三个二进制输入 $a,b,c$，总共只有 $2^3=8$ 种输入。表~\ref{l10:tab:parity} 给出一个完整的输入输出关系。只要列出全部输出为 1 的输入组合，就能构造对应的逻辑表达式和逻辑电路。

\begin{table}[H]
\centering
\caption{三个输入的奇偶性函数。}
\label{l10:tab:parity}
\begin{tabular}{cccc@{\qquad}l}
\toprule
$a$ & $b$ & $c$ & 输出 & 输入中 1 的个数\\
\midrule
0&0&0&0&0\\
0&0&1&1&1\\
0&1&0&1&1\\
0&1&1&0&2\\
1&0&0&1&1\\
1&0&1&0&2\\
1&1&0&0&2\\
1&1&1&1&3\\
\bottomrule
\end{tabular}
\end{table}

这张表还具有一个更紧凑的描述：输入中 1 的个数为奇数时输出 1，为偶数时输出 0。相应的布尔表达式是 $a\oplus b\oplus c$，其中 $\oplus$ 表示异或。\textbf{奇偶性就是能够解释全部观测的一个特征。}只有部分观测时，例如 $000\mapsto0$、$011\mapsto0$、$100\mapsto1$、$110\mapsto0$，还不能仅凭查表覆盖所有输入；能够提取这种规律，才有机会对未出现的组合给出答案。

\subsubsection{图像输入使查表规模呈指数增长}
一张 $32\times32$ 的图像包含 1024 个像素。按每像素 8 位计算，输入总共需要 8192 位，因此所有可能输入的数量是 $2^{8192}$。即使每种输入只对应一个类别，完整真值表仍然需要为这 $2^{8192}$ 种输入分别记录输出。

因为 $10^9\approx2^{30}$，十亿个不同输入覆盖的比例约为
\begin{equation}
\frac{2^{30}}{2^{8192}}=2^{-8162}.
\label{l10:eq:coverage}
\end{equation}
这个数量级说明，收集更多样本仍然无法接近穷举整个输入空间。学习过程需要利用多个样本共有的特征，把有限观测中的结构推广到未观测输入。小规模逻辑例子的奇偶性，正好说明了“从大量具体情况中提取简洁规律”这一作用。
